\chapter
[\CO{[RWT]} Rettungswesen und Einsatztaktik]
{\CC{[RWT]}\\Rettungswesen und Einsatztaktik}


\ChapterIntro
{%
Einführung ist die Strukturen des Rettungswesens und der
Einsatztaktik.

\noindent\SectionUnderConstruction
}
{%
\begin{description}
  \item[Maintainer:] Sebastian Gabriel
  \item[Autoren:] Diverse
  \item[Reviewer:] --
  \item[Version:] {\DocVersion} ({\DocVersionInfo})
  \item[SHA1:] {\DocGitCommit}
\end{description}

{\TxTeilkapitelAASS}
}




\section{Der Rettungs- und Krankentransportdienst}
\label{topic:rettungswesen-recht}

\Layout{2}{Krankentransportdienst}{%
Aufgabe eines \E{Krankentransportdienstes} ist es, Personen,
bei denen während des Transports eine Betreuung durch
Sanitäter medizinisch notwendig ist und die aus medizinischen
Gründen kein gewöhnliches Verkehrsmittel benützen können,
unter sachgerechter Betreuung mit geeigneten Transportmitteln
zu befördern.
}{}{}{}{}


\Layout{2}{Rettungsdienst}{%
Zu den Aufgaben eines \E{Rettungsdienstes} gehört es: 

\begin{enumerate}
  \item Personen, die eine
  erhebliche Gesundheitsstörung oder erhebliche Verletzung erlitten haben,
  \E{erste Hilfe} zu leisten, sie \E{transportfähig zu machen} und sie
  erforderlichenfalls \E{unter sachgerechter Betreuung} mit geeigneten
  Transportmitteln \E{in eine Krankenanstalt zu befördern} oder ärztlicher Hilfe
  zuzuführen;
  \item Personen wegen unmittelbarer
  Lebensgefahr sofortige erste \E{notärztliche} Hilfe zu leisten, die anders
  nicht gewährleistet ist;
  \item den Transport von Personen durchzuführen, bei denen
  \E{lebenswichtige Funktionen ständig überwacht oder aufrecht erhalten werden
  müssen};
  \item akute \E{Blut-, Blutprodukte- oder Organtransporte} durchzuführen;
  \item \E{Sanitätsdienste} zur Behandlung von akuten Erkrankungen oder
  Verletzungen bei \E{Veranstaltungen} mit dem hierfür erforderlichen Personal,
  den erforderlichen Einrichtungen und erforderlichen Transportmitteln bereit zu stellen;
  \item die
  Bevölkerung in erster Hilfe zu \E{schulen};
  \item im zivilen \E{Katastrophenschutz} mitzuwirken.
\end{enumerate}
}{}{}{}{}
















\Layout{2}{Kennzeichnung}{%
Die Verwendung der Bezeichnungen "`Rettungsdienst"' oder
"`Krankentransportdienst"' und von Übersetzungen dieser
Bezeichnungen in andere Sprachen in Geschäftspapieren,
Beschriftungen, Firmennamen oder Beschriftungen auf
Transportmitteln, die fälschlicherweise den Anschein
erwecken, dass es sich um einen Rettungs- oder
Krankentransportdienst handelt, ist zumeist durch
Landesgesetze verboten.

Die bei Rettungs- und Krankentransportdiensten tätigen
Personen haben im Dienst eine Kennzeichnung der Organisation
deutlich sichtbar zu tragen.
}{}{}{}{}


\Layout{2}{Der Ärztliche Leiter}{%
Rettungsdienste haben einen ärztlichen Leiter sowie einen
Stellvertreter zu bestellen. Der Ärztliche Leiter und deren
Stellvertreter in Abwesenheit des ärztlichen Leiters sind für
den  gesamten medizinischen Bereich des Rettungs- oder
Krankentransportdienstes verantwortlich.
}{}{}{}{}










\clearpage % TODO temp clearpage
\section{Kommunikationswege}
\label{topic:kommunikationswege}

\subsection{Die Leitstelle}
\label{topic:leitstelle}

\Layout{20}{\TxBeschreibung}{%
Eine Leitstelle leitet den
Betrieb einer oder mehrerer zugeordneten
Organisationseinheiten. Sie nimmt Informationen
(Notrufe, Meldungen der Einheiten, \ldots) entgegen, wertet
sie aus und koordiniert die einzelnen, ihr zugeordneten
Einheiten. 

Eine Leitstelle kann sowohl organisationsübergreifend (eine
Leitstelle ist für mehrere Organisationen zuständig,
z.\,B. für Rettungsdienst \E{und} Feuerwehr) eingerichtet
werden, es kann jedoch auch innerhalb einer Organisation
verschiedene Leitstellen, getrennt nach
Zuständigkeitsbereichen, geben (z.\,B. getrennte
Leitstellen für Rettungs- und Krankentransport,
Sonderleitstellen während Veranstaltungen, regionale
Trennung, \ldots)

Leitstellen können ortsfest eingerichtet werden, es
gibt jedoch auch mobile Leitstellen, welche bei besonderen
Ereignissen zum Einsatz kommen können.

{\TakeNotesMedium}
}
{}
{./images/gabriel-sebastian-aass/2010-06-25-DIF2010-IMG_2460_v2-quick00800-485.jpg}
{}
{GABS.}

\Layout{0}{Aufgaben}{%
Zu den typischen Aufgaben einer Rettungsdienstleitstelle
gehören die \EE{Annahme und Abfrage von Notrufmeldungen}
oder sonstigen Alarmierungen, sowie deren \EE{Erfassung} in
einem Einsatzleitsystem durch einen \E{Call Taker}. Die
Meldungen unterlaufen eine \EE{Einteilung und Priorisierung}
nach Art und Dringlichkeit der Meldung. Der Leitstelle
obliegt die Entscheidung, \E{welche Einsatzmittel}
eingesetzt werden, oft wird dazu nach einer Alarm- und
Ausrückordnung vorgegangen. Die \EE{Alarmierung} der
Einheiten erfolgt durch die Leitstelle über Sprech- oder
Datenfunk, bzw. andere geeignete
Hilfsmittel. Sie  \EE{übermittelt die Einsatzaufträge} und
unterstützt und koordiniert während des Einsatzes.

Die Einteilung, Alarmierung und Koordination der Einheiten
wird \E{Disposition} genannt und erfolgt durch
\T[Disponent]{Disponenten}.

In vielen Rettungsdienstbereichen ist die Leitstelle auch
für die Zuteilung von Krankenhausaufnahmekapazitäten
(\EE{Bettenzuteilung}) zuständig.
     
Um einen Überblick über einsatzbereite Fahrzeuge etc. zu erhalten, sorgt die
Leitstelle für eine \EE{Status- und Standort-Verwaltung}: Es
wird der jeweilige Status und -- wenn möglich -- der
Standort der Einheiten (einsatzbereit, Richtung
Berufungsort, \ldots) protokolliert und ausgewertet.

Eine sehr wichtige Aufgabe ist die \EE{Dokumentation}: Es
werden sowohl die Einsatzdaten wie z.\,B. Einsatzort,
Einsatzursache und Zeiten protokolliert, als auch besondere
Ereignisse. Telefongespräche mit der Leitstelle werden
i.\,d.\,R. aufgezeichnet und gespeichert. 
\A{Verweis Rechtsteil}

{\TakeNotesMedium} .
}{%
\begin{itemize}
  \item \E{Call Taking}
  \begin{itemize}
    \item Annahme \&\ Abfrage
  von Notrufen \&\ Alarmierungen 
  \item Erfassung im Einsatzleitsystem, Einteilung und Priorisierung nach Art
  und Dringlichkeit der Meldung, benötigtes Personal, etc. 
  \end{itemize} 
  \item \E{Disposition}
\begin{itemize}
    \item Welche Einsatzmittel
    \item Alarmierung der Einheiten
    \item Übermittlung der Aufträge
    \item Unterstützung, Koordination
    \item Evtl. Zuteilung von KH-Betten
    \item Status-/Standort-Verwaltung
  \end{itemize}
  \item Dokumentation
\end{itemize}
    
}{}{}{}


\subsection{Funk}
\index{Funk}

Als Funk(technik) bezeichnet man allgemein die
Methode, Informationen aller Art mit Hilfe
elektromagnetischer Wellen im Radiofrequenzbereich (Radiowellen) drahtlos zu
übertragen. 

\subsection{Sprechfunk}

Das drahtlose Übertragen von \E{gesprochenen} Informationen
wird als Sprechfunk bezeichnet. Dabei können
Handfunkgeräte, Mobilfunkgeräte oder Feststationen zum
Einsatz kommen.

Der klassische Funkbetrieb wird im
Wechselverkehr (\I[halbduplex]{Halbduplexbetrieb})
abgewickelt, d.\,h. es kann nur ein Teilnehmer gleichzeitig senden und belegt dabei
die Frequenz. \Z{Beim \I[vollduplex|Z]{Vollduplexbetrieb}
kann man gleichzeitig empfangen und senden (z.\,B. Telefon).} Um
einen geordneten Funkbetrieb zu ermöglichen, gibt es für
jeden organisierten Funkbereich Vorschriften (Protokolle),
wie der Funkverkehr abzuwickeln ist.


\Layout{2}{Funkname}{%
I.\,d.\,R. wird jede Funkstelle mit einem eigenen Funknamen
(Rufnahmen) bezeichnet und kann mit diesem gerufen werden.
Die Systematik der Bennenung kann je nach Land,
Organisation, Leitstellen-  und Funkbereich sehr unterschiedlich
sein. Die entsprechenden Protokolle sind zu beachten. 
}{}{}{}{}




\Fazit{Grundsätzlicher Ablauf im Sprechfunkverkehr:
\EE{Denken} -- \EE{drücken} -- \EE{kurz und bündig
sprechen}}

\Layout{2}{Beispiel}{%
Funkwagen 249 trifft am Berufungsort ein:

\begin{quote}
  \begin{labeling}[:]{ABK-Donaustadtx}
  \item[249] \Speech{"`Leitstelle für Zwoneunundvierzig."'}
  \item[Leitstelle] \Speech{"`Kommen Zwoneunundvierzig."'}
  \item[249] \Speech{"`Zwoneunundvierzig ist am BO eingetroffen."'}
  \item[Leitstelle] \Speech{"`Danke verstanden Zwoneunundvierzig."'}
\end{labeling}
\end{quote}
}{}{}{}{}



\Layout{2}{Beispiel}{%
Der RTW \E{ABK-Donaustadt} erhält über Sprechfunk einen
Einsatz:

\begin{quote}
  \begin{labeling}[:]{ABK-Donaustadtx}
  \item[Leitstelle] \Speech{"`ABK-Donaustadt für die
  Leitstelle."'}
  \item[ABK-Donaustadt] \Speech{"`ABK-Donaustadt hört."'}
  \item[Leitstelle] \Speech{"`Neuer Einsatz -- Zweiter
  Bezirk, Musterstraße -- Musterstraße
  Einhundertneunundfünzig, Stiege Zwo, dritter Stock, Tür
  dreiundvierzig -- mit sechsundzwanzig-Alpha-eins und
  zweihundertfünfundsechzigtausend-dreihundertsiebenundzwanzig."'}
  \item[ABK-Donaustadt]
  \Speech{"`ABK-Donaustadt hat verstanden und fährt
  Richtung BO."'}
\end{labeling}
\end{quote}
}{}{}{}{}



 


\Layout{57}{}{}{}
{%
\begin{tabular}{clll|clll|cll}
 & \TabHeading{ÖNORM} & \TabHeading{DIN}		& \TabHeading{NATO} &  & \TabHeading{ÖNORM} & \TabHeading{DIN} & \TabHeading{NATO} 
\\\addlinespace\midrule\addlinespace
\TabSub{A} & Anton 		& $\hookleftarrow$	& Alpha 	& \TabSub{O} & Otto 		& $\hookleftarrow$	& Oscar 	& \TabSub{0} & Null 	& Zero\\\addlinespace
\TabSub{Ä} & Ärger 		& $\hookleftarrow$	& --- 		& \TabSub{Ö} & Österreich 	& Ökonom	& --- 		& \TabSub{1} & Eins 	& One\\\addlinespace
\TabSub{B} & Berta 		& $\hookleftarrow$	& Bravo 	& \TabSub{P} & Paula 		& $\hookleftarrow$	& Papa 		& \TabSub{2} & Zwo 		& Two\\\addlinespace
\TabSub{C} & Cäsar 		& $\hookleftarrow$	& Charlie 	& \TabSub{Q} & Quelle 		& $\hookleftarrow$	& Quebec 	& \TabSub{3} & Drei 	& Three\\\addlinespace
\TabSub{Ch}& --- 		& Charlotte	& --- 		& \TabSub{R} & Richard 		& $\hookleftarrow$	& Romeo 	& \TabSub{4} & Vier 	& Fower\\\addlinespace
\TabSub{D} & Dora 		& $\hookleftarrow$	& Delta 	& \TabSub{S} & Siegfried 	& Samuel	& Sierra 	& \TabSub{5} & Fünf 	& Five\\\addlinespace
\TabSub{E} & Emil 		& $\hookleftarrow$	& Echo 		& \TabSub{Sch} & --- 		& Schule	& ---		& \TabSub{6} & Sechs 	& Six\\\addlinespace
\TabSub{F} & Friedrich 	& $\hookleftarrow$	& Foxtrott 	& \TabSub{ß} & Scharfes s	& Eszett	& --- 		& \TabSub{7} & Sieben 	& Seven\\\addlinespace
\TabSub{G} & Gustav 	& $\hookleftarrow$	& Golf 		& \TabSub{T} & Theodor 		& $\hookleftarrow$	& Tango 	& \TabSub{8} & Acht 	& Eight\\\addlinespace
\TabSub{H} & Heinrich 	& $\hookleftarrow$	& Hotel 	& \TabSub{U} & Ulrich 		& $\hookleftarrow$	& Uniform 	& \TabSub{9} & Neun 	& Niner\\\addlinespace
\TabSub{I} & Ida 		& $\hookleftarrow$	& India 	& \TabSub{Ü} & Übel 		& Übermut	& --- 		& \TabSub{./,} & --- 	& Decimal\\\addlinespace
\TabSub{J} & Julius 	& $\hookleftarrow$	& Juliett 	& \TabSub{V} & Viktor 		& $\hookleftarrow$	& Victor 	& \TabSub{.} & --- 		& Stop\\\addlinespace
\TabSub{K} & Konrad 	& Kaufmann			& Kilo 		& \TabSub{W} & Wilhelm 		& $\hookleftarrow$	& Whiskey	& \\\addlinespace
\TabSub{L} & Ludwig 	& $\hookleftarrow$	& Lima 		& \TabSub{X} & Xaver 		& Xanthippe	& X-ray		& \\\addlinespace
\TabSub{M} & Martha 	& $\hookleftarrow$	& Mike 		& \TabSub{Y} & Ypsilon 		& $\hookleftarrow$	& Yankee\\\addlinespace
\TabSub{N} & Nordpol 	& $\hookleftarrow$	& November	& \TabSub{Z} & Zürich 		& Zacharias	& Zulu\\\addlinespace
\\\addlinespace\bottomrule
\end{tabular}
}
{Buchstabieralphabete nach ÖNORM A 1081, DIN 5009 und NATO}
{}







{\TakeNotesMedium}


\subsection{Datenfunk}

Unter dem Begriff \I{Datenfunk} versteht man die
\E{automatisierte}, drahtlose Übertragung von codierten
\EE{Daten} mittels analoger oder digitaler Funktechnik.
I.\,d.\,R. erfolgt dies mittels EDV-gestützter Systeme,
welche sich um die Codierung und anwendungsgerechte
Darstellung (Monitor, Fahrzeugdisplays, etc.) der Daten
kümmern. Der Vorteil des Datenfunkes ist, dass Informationen
nur einmal eingegeben werden müssen, und von den Systemen
der Empfangsstationen (Fahrzeugdisplays, elektronische
Dokumentationssysteme, etc.) automatisiert weiterverwendet
werden können. Am Markt haben sich eine Vielzahl von
Produkten fest etabliert und lösen den Sprechfunk als
Kommunikationsmedium immer mehr ab.

{\TakeNotesMedium}

\subsection{Telefon}

Das Telefon oder Mobiltelefon hat sich in der täglichen
Routine etabliert. Nachteil dieses Systems ist, dass man auf
die zivile Infrastruktur angewiesen ist. Ist diese z.\,B.
in Folge von Überlastung gestört, ist dieser
Kommunikationsweg nicht mehr zuverlässig nutzbar und man
muss mit den Einsatzorganisationen vorbehaltenen
Kommunikationswegen auskommen.

{\TakeNotesMedium}

\section{Einsatzmittel}
\label{topic:einsatzmittel}

Im deutschsprachigen Raum hat sich die Einteilung in die
Typen 

\begin{itemize}
  \item \EE{Krankentransportwagen} (\EE{KTW}), 
  \item \EE{Rettungstransportwagen}
(\EE{RTW}), 
\item \EE{Notarztwagen} (\EE{NAW}), 
\item \EE{Notarzteinsatzfahrzeug}
(\EE{NEF}) und 
\item \EE{Notarzthubschrauber} (\EE{NAH}) (in Deutschland:
Rettungstransporthubschrauber, RTH) 
\end{itemize}
eingebürgert.

\Layout{2}{Krankentransportwagen -- KTW}
{%
Der \Abk{KTW} ist für den Transport, die Erstversorgung und
die Überwachung von Patienten, welche \E{keine
Notfallpatienten} sind und keine intensive medizinische
Betreuung benötigen, konstruiert und ausgerüstet.

Das Haupteinsatzgebiet ist der \EE{Krankentransport}, daneben
wird der \Abk{KTW} auch zum Transport von erkrankten,
unkritischen Patienten im Rahmen des Rettungsdienstes
eingesetzt.

Im \EE{Ausnahmefall}, wenn z.\,B. kein geeigneteres Fahrzeug
in angemessener Zeit verfügbar ist, oder im Großschadensfall,
können auch Notfallpatienten mit einem \Abk{KTW}
transportiert werden.

Ein \Abk{KTW} ist mit mindestestens 2 Fachpersonen besetzt
(Mindestqualifikation: 2 RS).

{\TakeNotesMedium}
% \blindtext

}{

}
{./icons/under-construction.png}
{}
{}

\Layout{60}{Rettungstransportwagen -- RTW}
{%
Der \Abk{RTW} ist für den Transport, die erweiterte
Behandlung und die Überwachung von Patienten -- insbesonders
Notfallpatienten -- konstruiert und ausgerüstet. Er wird in
der Notfallrettung und bei Überstellungen von Patienten,
welche eine erweiterte Behandlung benötigen, eingesetzt.

Steigt ein Notarzt zu, sollte der \Abk{RTW} weitgehend
gleichwertig zu einem \Abk{NAW} sein\ZFootnote{Aufgrund
diverser Umstände und Vorgaben kann die Forderung, dass ein
\Abk{RTW} einem \Abk{NAW} ohne Notarzt entspricht, nicht
immer uneingeschränkt umgesetzt werden. Probleme ergeben sich
insbesonders durch das Suchtmittelgesetz und besonders teure
Geräte oder Medikamente. In jedem Fall soll aber eine
"`Umrüstung"' vom \Abk{RTW} hin zum \Abk{NAW} während des
Einsatzes mit \E{wenigen Handgriffen} möglich sein.}.
Zusammen mit dem \Abk{NEF} führt der \Abk{RTW}
notarztpflichtige Einsätze durch.

Ein \Abk{RTW} ist mit mindestens 2, oft auch 3 Fachpersonen besetzt
(Mindestqualifikation: 1 \Abk{NFS}, \VonBis{1}{2} \Abk{RS}).

{\TakeNotesMedium}
% \blindtext
}{

}
{./images/teaser/rtw_ma70_asb921.jpg}
{RTWs/NAWs der Wiener Berufsrettung (li.) und des ASB (re.).}
{Gabriel.}


\Layout{60}{Notarztwagen -- NAW}
{%
Der \Abk{NAW} ist für den Transport, die erweiterte
\E{ärztiche} Behandlung und die Überwachung von Patienten --
insbesonders Notfallpatienten -- konstruiert und ausgerüstet.
Er wird in der Notfallrettung und bei Überstellungen von
Patienten, welche eine ärztliche Behandlung oder Überwachung
benötigen, eingesetzt.
% Notfallrettung \& Überstellungen.

Der \Abk{NAW} ist mit mindestens 3, oft auch 4 Fachpersonen besetzt
(Mindestqualifikation: 1 \Abk{NA}, 1 \Abk{NFS}, \VonBis{1}{2}
\Abk{RS}).

{\TakeNotesMedium}

% \blindtext

}{

}
{./images/teaser/rtw_ma70_asb921.jpg}
{RTWs/NAWs der Wiener Berufsrettung (li.) und des ASB (re.).}
{Gabriel.}


\Layout{60}{Notarzteinsatzfahrzeug -- NEF}
{%
Ein \Abk{NEF} ist ein Zubringerfahrzeug, welches den Notarzt
zu einem Einsatzort bringt. Es führt die für eine ärztliche
Behandlung notwendige Ausstattung mit sich. Ein \Abk{NEF}
kann zwar grundsätzlich ohne andere Rettungsmittel Patienten
versorgen, ein Transport muss jedoch mit einem anderen
Fahrzeug (\Abk{RTW}, \Abk{KTW}, \ldots) erfolgen.

Das \Abk{NEF} ist mit mindestens 2 Personen besetzt,
(Mindestqualifikation: 1 \Abk{NA}, 1 \Abk{NFS} (im
Ausnahmefall auch \Abk{RS})).

\Fazit{Ein NEF transportiert keine Patienten.} 



{\TakeNotesMedium}

% \blindtext
}{

}
{./images/teaser/nef-passat-kontur.jpg}
{NEF der Wiener Berufsrettung (altes Modell).}
{Gabriel.}


\Layout{60}{Notarzthubschrauber -- NAH}
{%
Notarzthubschrauber werden in der Notfallrettung, für
Überstellungen und z.\,T. bei Bergungen eingesetzt.

Vorteile sind u.\,a. ein relativ erschütterungsfreier
Transport und kürzere Transportzeiten auf längeren Distanzen.
Der \Abk{NAH} hat aber auch eine Reihe von \EE{Nachteilen}:

\begin{itemize}
  \item \EE{Landefläche} Der Hubschrauber benötigt einen
  geeigneten Landeplatz, i.\,d.\,R. mind.
  25\,$\times$\,25\Meter*, es dürfen keine Oberleitungen
  o.\,ä. den gefahrlosen Anflug stören. Ist in der Nähe des
  Einsatzortes keine geeignete Landefläche verfügbar muss der
  Patient erst mittels eines Fahrzeuges zu einem
  entsprechenden Platz transportiert werden, dabei geht Zeit
  verloren.
  \item \EE{Wetter} Geeignetes Flugwetter ist Voraussetzung.
  \item \EE{Tageszeit} Nur wenige Hubschrauber sind
  nachtflugtauglich, außerdem verbietet sich eine Landung
  abseits gesicherter Landeplätz in der Nacht.
  \item \EE{Kosten} Sowohl die Vorhaltung als auch der Einsatz
  eines Hubschraubers ist immens teuer.
  \item \EE{Platz} Je nach Typ ist im Hubschrauber aufgrund des
  beschränkten Platzangebotes selber
eine Behandlung und Betreuung oft nur eingeschränkt möglich.
\end{itemize}

Die Vor- und Nachteile müssen abgewogen werden. So kann es
bei dem gleichen Patienten mit der gleichen Diagnose klug
sein mit dem \Abk{NAH} zu transportieren, wenn der Notfall
während der Hauptverkehrszeit eingetreten ist, wohingegen zu
einer anderen Zeit evtl. ein bodengebundener Transport
schneller und vorteilhafter ist.



Die Besetzung ist uneinheitlich, zwingend erforderlich sind
jedoch mind. 1 Pilot und 1 Notarzt.

{\TakeNotesMedium}
% \blindtext
}{

}
{./images/gabriel-sebastian-aass/Landeanflug-OE-XEC_IMG_1912_1.jpg}
{Ein Hubschrauber der Christophorus-Flugrettungsflotte im Landeanflug an das AKH
Wien.} {GABS.}


\Layout{2}{Weitere Fahrzeugtypen}
{
PKW, N-KTW, B-KTW, SEW, MLS, ITW

{\TakeNotesLarge}
% \blindtext
}{

}
{./icons/under-construction.png}
{}
{}

% \begin{WidePar}

\Layout{57}{}
{}
{}
{\small%
\begin{tabularx}{\linewidth}{lXXXXX}

&
\TabHeading{KTW}
&
\TabHeading{RTW}
&
\TabHeading{NAW}
&
\TabHeading{NEF}
&
\TabHeading{NAH}
\\\addlinespace\midrule\addlinespace
\TabSub{Name}
&
Kranken\-transport\-wagen
&
Rettungs\-transport\-wagen
&
Notarzt\-wagen
&
Notarzt\-einsatz\-fahrzeug
&
Notarzt\-hubschrauber
\\\addlinespace
\TabSub{Besatzung}
&
2 RS
&
1 NFS, \VonBis{1}{2} RS
&
1 NA, 1 NFS, \VonBis{1}{2} RS
&
1 NFS, 1 NA
&
1 Pilot, 1 NFS, 1 NA
\\\addlinespace
\TabSub{Funktion}
&	
Krankentransport.

Transport von erkrankten Patienten, die keine intensiven med. Maßnahmen
benötigen. 
&
Notfallrettung \& Überstellungen. 

Versorgung und Transport von erkrankten Patienten und Notfallpatienten. 
&
Notfallrettung \& Überstellungen. 

Versorgung und Transport von Notfallpatienten. 
&
Notfallrettung \& Überstellungen. 

Versorgung von erkrankten Notfallpatienten, \E{kein Transport}. Zubringer von
Notarzt. 
&
Notfallrettung \& Überstellungen. 

Versorgung und Transport von Notfallpatienten. 
\\\addlinespace
\TabSub{Anmerkungen}

&

&
Steigt ein Notarzt zu, sollte der RTW weitgehend gleichwertig zu einem NAW sein. 
&

&

&
In Deutschland: "`RTH"'.
\\\addlinespace\bottomrule
\end{tabularx}
}
{Übersicht der wichtigsten Einsatzmitteltypen in Österreich}
{}


% \end{WidePar}

\section{Ablauf eines Einsatzes}
\label{topic:einsatzarten}

\subsection{Primärtransporte}

{\TakeNotesLarge}

\subsection{Sekundärtransport}

{\TakeNotesLarge}

\section{Regionales Rettungswesen}

\Layout{2}
{\TxQuerverweise}
{~\par\noindent\includegraphics[angle=270,width=\linewidth]{./images/leitsystem/Leitschema-PAM_SAN.pdf}}
{\begin{itemize}
  \item \EE{Regionele Besonderheiten / Regionales
  Rettungswesen}:
  \FullPrettyRef{topic:regionales-rettungswesen}
\end{itemize}}
{}
{}
{}

